\documentclass[handout]{beamer}
\mode<presentation>
\usepackage{xcolor}
\usepackage[toc]{multitoc}
\usepackage{color}
\usepackage{graphicx}
\graphicspath{{../buku_ajar/buku_ajar_logika_matematika/figures/}}
\usepackage{tikz}
\usetikzlibrary{shapes,arrows,positioning,calc}
\usepackage{listings}
\usepackage{lmodern}
\usepackage{amsmath}
\usetheme[secheader]{Boadilla}
\usecolortheme{whale}
\setbeamertemplate{caption}[numbered]
\setbeamertemplate{bibliography item}[text]
\setbeamertemplate{navigation symbols}{}

\newcommand{\logic}[1]{\textbf{\textit{#1}}}

\theoremstyle{definition}
\newtheorem{definisi}[theorem]{Definisi}
\theoremstyle{example}
\newtheorem{contoh}[theorem]{Contoh}

\renewcommand*{\multicolumntoc}{2}

\setbeamertemplate{footline}
{
  \leavevmode%
  \hbox{%
  \begin{beamercolorbox}[wd=.333333\paperwidth,ht=2.25ex,dp=1ex,center]{author in head/foot}%
    \usebeamerfont{author in head/foot}\insertshortauthor%
  \end{beamercolorbox}%
  \begin{beamercolorbox}[wd=.433333\paperwidth,ht=2.25ex,dp=1ex,center]{title in head/foot}%
    \usebeamerfont{title in head/foot}\insertshorttitle
  \end{beamercolorbox}%
  \begin{beamercolorbox}[wd=.233333\paperwidth,ht=2.25ex,dp=1ex,right]{date in head/foot}%
    \usebeamerfont{date in head/foot}\insertshortdate{}\hspace*{2em} \insertframenumber{} / \inserttotalframenumber\hspace*{2ex} 
  \end{beamercolorbox}}%
  \vskip0pt%
}

\subtitle[Logika Matematika]{Logika Matematika}
\title[Peta Karnaugh]{Penyederhanaan Fungsi Logika: Aljabar Boolean dan Peta Karnaugh}
\author{Cecep Suwanda, S.Si., M.Kom.}
\date{Maret 2025}
\institute{Universitas Bale Bandung}

\begin{document}

\maketitle

\section*{Daftar Isi}
\begin{frame}
\frametitle{Daftar Isi}
\tableofcontents
\end{frame}

% ============================================================
% SECTION 10-1: Motivasi
% ============================================================
\section{Motivasi dan K-Map}

\begin{frame}
\frametitle{Keterbatasan Metode Aljabar}
\begin{itemize}
  \item Tidak ada algoritma deterministik yang menjamin solusi optimal
  \item Keberhasilan bergantung pada intuisi dan pengalaman
  \item Untuk 5+ variabel: manipulasi rumit dan rentan kesalahan
\end{itemize}
\begin{block}{Peta Karnaugh (Maurice Karnaugh, 1953)}
Metode visual berbasis pengenalan pola pada grid 2D. Menggunakan Kode Gray. Efektif hingga 4 variabel; 5--6 variabel dengan penyesuaian. Untuk $>$6 variabel: Quine--McCluskey.
\end{block}
\end{frame}

\begin{frame}
\frametitle{Prinsip Kerja K-Map}
\begin{enumerate}
  \item Kombinasi input dipetakan ke sel grid 2D (setiap sel = 1 minterm)
  \item Tata letak: sel bersebelahan hanya berbeda 1 variabel (Kode Gray)
  \item Dua sel bersebelahan bernilai 1 $\Rightarrow$ variabel pembeda dieliminasi (hukum komplemen $x+x'=1$)
  \item Semakin besar kelompok $\Rightarrow$ semakin banyak variabel dieliminasi
\end{enumerate}
\end{frame}

% ============================================================
% SECTION 10-2: Kode Gray
% ============================================================
\section{Kode Gray}

\begin{frame}
\frametitle{Definisi Kode Gray}
\begin{definisi}[Kode Gray]
Sistem pengkodean biner di mana dua bilangan berurutan hanya berbeda pada tepat satu digit. Properti: \logic{jarak Hamming satu}.
\end{definisi}
Urutan label K-Map: $00, 01, 11, 10$ (bukan biner standar $00, 01, 10, 11$). Sifat siklis: baris terakhir dan pertama juga berbeda 1 bit $\Rightarrow$ \textbf{wrap-around}.
\end{frame}

\begin{frame}
\frametitle{Konstruksi Rekursif}
\begin{enumerate}
  \item Mulai: 1-bit $\{0, 1\}$
  \item Refleksikan urutan; tambah prefiks 0 pada asli, 1 pada refleksi
  \item Gabungkan
\end{enumerate}
2-bit: $00, 01, 11, 10$\\
3-bit: $000, 001, 011, 010, 110, 111, 101, 100$
\end{frame}

% ============================================================
% SECTION 10-3: K-Map 2 Variabel
% ============================================================
\section{K-Map 2 dan 3 Variabel}

\begin{frame}
\frametitle{Aturan Pengelompokan}
\begin{enumerate}
  \item Ukuran kelompok: pangkat 2 ($1, 2, 4, 8, \ldots$)
  \item Bentuk persegi/persegi panjang (bukan diagonal/L)
  \item Maksimalkan ukuran kelompok
  \item Minimalkan jumlah kelompok
  \item Semua sel 1 harus tercakup
  \item Tumpang tindih diperbolehkan
\end{enumerate}
\end{frame}

\begin{frame}
\frametitle{Eliminasi Variabel}
\begin{center}
\small
\begin{tabular}{|c|c|l|}
\hline
\textbf{Ukuran} & \textbf{Eliminasi} & \textbf{Suku} \\ \hline
1 sel & 0 var & Minterm lengkap \\ \hline
2 sel & 1 var & $(n-1)$ literal \\ \hline
4 sel & 2 var & $(n-2)$ literal \\ \hline
$2^k$ sel & $k$ var & $(n-k)$ literal \\ \hline
\end{tabular}
\end{center}
\end{frame}

\begin{frame}
\frametitle{Contoh K-Map 2 Var: $F = \sum m(0,1,3)$}
\begin{center}
\small
\begin{tabular}{|c|c|c|}
\hline
 & $Y=0$ & $Y=1$ \\ \hline
$X=0$ & 1 ($m_0$) & 1 ($m_1$) \\ \hline
$X=1$ & 0 & 1 ($m_3$) \\ \hline
\end{tabular}
\end{center}
Kelompok 1: $\{m_0, m_1\}$ $\Rightarrow$ $X'$\\
Kelompok 2: $\{m_1, m_3\}$ $\Rightarrow$ $Y$\\
\textbf{Hasil}: $F = X' + Y$
\end{frame}

% ============================================================
% SECTION 10-4: K-Map 3 Variabel
% ============================================================
\begin{frame}
\frametitle{K-Map 3 Variabel}
Grid $2 \times 4$. Kolom: $00, 01, 11, 10$ (Gray). \textbf{Wrap-around}: kolom pertama $\leftrightarrow$ kolom terakhir.
\begin{block}{Contoh: $F = \sum m(1,3,5,7)$}
Keempat sel 1 di kolom $Z=1$ $\Rightarrow$ $X,Y$ berubah $\Rightarrow$ $F = Z$
\end{block}
\end{frame}

\begin{frame}
\frametitle{Contoh Wrap-Around: $F = \sum m(0,2,4,5)$}
$\{m_0, m_2\}$: wrap kolom $00 \leftrightarrow 10$, $X'Z'$ konstan $\Rightarrow$ $X'Z'$\\
$\{m_4, m_5\}$: $XY'$ konstan $\Rightarrow$ $XY'$\\
\textbf{Hasil}: $F = X'Z' + XY'$
\end{frame}

% ============================================================
% SECTION 10-5: K-Map 4 Variabel
% ============================================================
\section{K-Map 4 Variabel}

\begin{frame}
\frametitle{K-Map 4 Variabel}
Grid $4 \times 4$; baris dan kolom mengikuti Gray. Wrap pada kedua sumbu. Empat sudut ($m_0, m_2, m_8, m_{10}$) bersebelahan logis.
\begin{center}
\small
\begin{tabular}{|c|c|l|}
\hline
\textbf{Ukuran} & \textbf{Eliminasi} & \textbf{Contoh} \\ \hline
4 sel & 2 var & $A'B'$, $C'D$ \\ \hline
8 sel & 3 var & $A'$, $D$ \\ \hline
16 sel & 4 var & $1$ (tautologi) \\ \hline
\end{tabular}
\end{center}
\end{frame}

\begin{frame}
\frametitle{Contoh: Empat Sudut}
$F = \sum m(0,2,8,10)$: 4 sel di pojok $\Rightarrow$ $C'D'$ (double wrap)\\
$F = \sum m(0,2,5,7,8,10,13,15)$: pojok $C'D'$ + tengah $CD$ $\Rightarrow$ $F = C'D' + CD$ (XNOR)
\end{frame}

% ============================================================
% SECTION 10-6: Don't Care
% ============================================================
\section{Don't Care dan Prime Implicant}

\begin{frame}
\frametitle{Kondisi Don't Care}
\begin{definisi}[Don't Care]
Kombinasi input yang tidak memiliki spesifikasi output. Dinotasikan $d$ atau $\times$. Perancang bebas pilih 0 atau 1 untuk optimasi.
\end{definisi}
Contoh: BCD menggunakan 4 bit untuk 0--9; kombinasi 10--15 = don't care.
\end{frame}

\begin{frame}
\frametitle{Prime Implicant dan Essential}
\begin{definisi}[Prime Implicant]
Kelompok sel 1 yang \textbf{maksimal}: tidak dapat diperbesar tanpa menyertakan sel 0.
\end{definisi}
\begin{definisi}[Essential Prime Implicant]
Prime implicant yang mencakup minimal satu sel 1 yang \textbf{tidak dicakup} prime implicant lain. Wajib disertakan.
\end{definisi}
Prosedur: identifikasi prime implicant $\rightarrow$ tentukan essential $\rightarrow$ pilih tambahan untuk sisa minterm.
\end{frame}

% ============================================================
% SECTION 10-7: Quine-McCluskey
% ============================================================
\section{Quine--McCluskey}

\begin{frame}
\frametitle{Keterbatasan K-Map}
\begin{center}
\small
\begin{tabular}{|c|c|l|}
\hline
\textbf{Var} & \textbf{Sel} & \textbf{Keterangan} \\ \hline
2--4 & 4--16 & Praktis \\ \hline
5 & 32 & Dua peta 4-var \\ \hline
6 & 64 & Empat peta \\ \hline
7+ & 128+ & Tidak praktis; perlu algoritma \\ \hline
\end{tabular}
\end{center}
\end{frame}

\begin{frame}
\frametitle{Algoritma Quine--McCluskey}
\begin{definisi}[Quine--McCluskey]
Prosedur algoritmik (tabulasi) untuk menemukan semua prime implicant dan memilih himpunan minimum. Quine (1952), McCluskey (1956).
\end{definisi}
Tahap 1: gabungkan minterm yang berbeda 1 bit, ulangi hingga selesai. Tahap 2: pilih himpunan minimal dengan tabel cakupan.
\end{frame}

\begin{frame}
\frametitle{K-Map vs Quine--McCluskey}
\begin{center}
\small
\begin{tabular}{|l|l|l|}
\hline
\textbf{Aspek} & \textbf{K-Map} & \textbf{Q--M} \\ \hline
Pendekatan & Visual & Algoritmik \\ \hline
Variabel efektif & 2--6 & Tidak terbatas \\ \hline
Optimalitas & Bergantung pengguna & Dijamin minimal \\ \hline
Otomatisasi & Kurang & Mendukung \\ \hline
\end{tabular}
\end{center}
\end{frame}

% ============================================================
% RANGKUMAN
% ============================================================
\section{Rangkuman}

\begin{frame}
\frametitle{Rangkuman}
\begin{itemize}
  \item K-Map: metode visual; label dengan Kode Gray; sel bersebelahan = berbeda 1 var
  \item Aturan: kelompok pangkat 2, persegi panjang, wrap-around
  \item Don't care: fleksibel 0/1 untuk optimasi
  \item Prime implicant maksimal; essential wajib disertakan
  \item Quine--McCluskey: algoritmik, skala ke banyak variabel
\end{itemize}
\end{frame}

\begin{frame}
\frametitle{Kata Kunci}
\begin{center}
\small
Peta Karnaugh $\bullet$ Kode Gray $\bullet$ wrap-around $\bullet$ pengelompokan\\
don't care $\bullet$ prime implicant $\bullet$ essential prime implicant $\bullet$ Quine--McCluskey
\end{center}
\end{frame}

\end{document}
